\documentclass[11pt,a4paper]{article}

% ============== ENCODING AND FONTS ==============
\usepackage[utf8]{inputenc}
\usepackage[T1]{fontenc}
\usepackage{lmodern}
\usepackage{microtype}

% ============== PAGE LAYOUT ==============
\usepackage[margin=1in]{geometry}
\usepackage{setspace}

% ============== MATHEMATICS ==============
\usepackage{amsmath,amssymb,amsfonts,amsthm,bm}
\usepackage{siunitx}
\usepackage{physics}     % for \dd, \dv, etc.; see clash fix below

% Fix siunitx + physics name clash for \qty
\AtBeginDocument{\RenewCommandCopy\qty\SI}

% ============== GRAPHICS AND TABLES ==============
\usepackage{graphicx}
\usepackage{booktabs}
\usepackage{float}

% ============== LISTS AND ENUMERATIONS ==============
\usepackage{enumitem}

% ============== BIBLIOGRAPHY ==============
% Using embedded bibliography; BibTeX integration can be added in future revisions

% ============== HYPERLINKS (load last) ==============
\usepackage[colorlinks=true,linkcolor=blue,citecolor=blue,urlcolor=blue]{hyperref}

% ============== CUSTOM MACROS ==============
\newcommand{\dx}{\mathrm{d}}
\newcommand{\Geff}{G_{\mathrm{eff}}}
\newcommand{\ac}{a_{\mathrm{c}}}
\newcommand{\zc}{z_{\mathrm{c}}}
\newcommand{\rs}{r_{\mathrm{s}}}
\newcommand{\cs}{c_{\mathrm{s}}}
\newcommand{\mpl}{M_{\mathrm{Pl}}}
\newcommand{\Hrh}{H_{\mathrm{rh}}}
\newcommand{\Trh}{T_{\mathrm{rh}}}
\newcommand{\sigmamech}{\sigma_{\mathrm{mech}}}
\newcommand{\fmech}{f_{\mathrm{mech}}}
\newcommand{\umech}{u_{\mathrm{mech}}}
\newcommand{\rhocrit}{\rho_{\mathrm{crit}}}
\newcommand{\rhostd}{\rho_{\mathrm{std}}}
\newcommand{\Zpf}{Z_{\mathrm{pf}}}
\newcommand{\Zcond}{Z_{\mathrm{c}}}
\newcommand{\uf}{u_{\mathrm{f}}}

% Theorem-like environments for postulates, propositions, etc.
\newtheorem{postulate}{Postulate}
\newtheorem{proposition}{Proposition}
\theoremstyle{definition}
\newtheorem{definition}{Definition}

% TODO markers for collaborative work
\newcommand{\TODO}[1]{\textcolor{red}{\textbf{[TODO: #1]}}}
\newcommand{\MATHAGENT}{\textcolor{orange}{\textbf{[MATH-AGENT]}}}
\newcommand{\PHYSAGENT}{\textcolor{teal}{\textbf{[PHYS-AGENT]}}}
\newcommand{\GRAPHAGENT}{\textcolor{purple}{\textbf{[GRAPHIC-AGENT]}}}

% ============== TITLE ==============
\title{\textbf{The Observable Universe as a Condensate Expanding into a Universal Protofluid}}
\author{Rob Skavberg\thanks{Independent Researcher, Calgary, Canada. Contact: \texttt{rob.skavberg@example.org}}}
\date{\today \\ \textit{Draft Version 0.1 --- For Team Collaboration}}

\begin{document}
\maketitle
\doublespacing

% ============== ABSTRACT ==============
\begin{abstract}
We present a foundational framework in which the \emph{observable Universe} is interpreted as a quantum condensate expanding into a universal protofluid. In this picture, the Big Bang corresponds to the nucleation and subsequent growth of a phase boundary, rather than an initial singularity. Gravitational response emerges from the medium's effective stiffness---a mechanical-geometric linkage elaborated in prior work \cite{SkavbergMechG2025,SkavbergMechGR2025,SkavbergPlanckField2025}---while boundary-layer impedance mismatch mediates reflection, refraction, and transient energy reservoirs that can briefly modify the early expansion history in a testable epoch.

A central contribution of this paper is a rigorous \emph{Validation of Foundational Proofs} section, which establishes quantitative pass/fail criteria against well-established pillars: classical General Relativity tests (Solar System, binary pulsars), gravitational-wave propagation (GW170817 luminal bound), Big-Bang Nucleosynthesis (BBN), the Cosmic Microwave Background (CMB/Planck 2018), Baryon Acoustic Oscillations (BAO), Type Ia supernovae (SNe), distance ladders, and large-scale structure (LSS). These guardrails ensure the framework reduces to standard physics in all tested regimes, with deviations permitted only in a narrow, early-time window that can address persistent cosmological tensions (e.g., the Hubble constant discrepancy).

We conclude with a sequenced research roadmap comprising seven follow-on papers designed to stress-test the framework empirically, from early-Universe baseline fits and BBN compatibility, through gravitational-wave propagation and boundary-layer kinematics, to structure formation and targeted null tests. This draft serves as the architectural foundation for collaborative theoretical and computational work by the research team.
\end{abstract}

% ============== TABLE OF CONTENTS ==============
\tableofcontents
\newpage

% ============== SECTION 1: INTRODUCTION ==============
\section{Introduction}
\label{sec:intro}

The standard $\Lambda$CDM cosmological model has achieved remarkable empirical success, fitting a wide array of observations---from the cosmic microwave background (CMB) \cite{Planck2018} to large-scale structure (LSS), baryon acoustic oscillations (BAO), and Type Ia supernovae (SNe) \cite{RiessReview2024}---with a minimal parameter set. Yet persistent tensions have emerged: most notably, the $\sim$5$\sigma$ discrepancy between early-Universe (CMB-inferred) and late-Universe (distance-ladder) determinations of the Hubble constant $H_0$ \cite{RiessReview2024,EDEReview2023}, and the $S_8$ tension in measurements of matter clustering amplitude \cite{PlanckS8,DES2021}. These anomalies have motivated exploration of extensions to $\Lambda$CDM, including early dark energy (EDE) models \cite{PoulinReview2023,EDEReview2023}, modified gravity scenarios, and new physics in the dark sector.

Parallel to these observational puzzles are long-standing conceptual questions: What is the microphysical origin of gravitational response? Why does the Einstein--Hilbert action describe spacetime curvature so effectively? What are the initial conditions for the Big Bang, and is the initial singularity physically real or an artifact of the classical approximation? Inspired by induced-gravity frameworks \cite{Sakharov,Vassilevich} and analogue-gravity systems \cite{AnalogueGravityReview}, a growing body of work explores the possibility that spacetime geometry and gravitational dynamics are emergent phenomena arising from an underlying condensed-matter-like substrate.

\subsection{Motivation: A Condensate-Based Universe}
\label{sec:motivation}

This paper advances a bold organizing principle: the \emph{observable Universe is a coherent quantum condensate} that nucleated and expanded into a pre-existing \emph{universal protofluid}. In this picture:
\begin{itemize}[leftmargin=1.8em]
  \item The \textbf{Big Bang} is reinterpreted as the \emph{nucleation event} of a phase transition, initiating the formation of a codimension-one conversion boundary separating the protofluid phase from the condensate phase.
  \item \textbf{Cosmic expansion} corresponds to the outward propagation of this conversion boundary, with the condensate interior expanding as the boundary advances.
  \item \textbf{Gravitational response} emerges from the condensate's effective stiffness (bulk modulus, phase coherence, and quantum pressure), as detailed in prior work on the mechanical route to $G$ \cite{SkavbergMechG2025}, the mechanical extension of General Relativity \cite{SkavbergMechGR2025}, and Planck-field resolution of the Hulse--Taylor binary \cite{SkavbergPlanckField2025}.
  \item \textbf{Boundary-layer physics} mediates impedance mismatch between protofluid and condensate, leading to reflection, refraction, and \emph{transient energy accumulation} that can briefly alter the early expansion rate $H(a)$, offering a testable mechanism to address the $H_0$ tension.
\end{itemize}

\TODO{PHYS-AGENT: Expand this subsection with additional motivation from analogue gravity literature, condensed-matter analogues (BEC, superfluid He, etc.), and recent CMB tension reviews.}

\subsection{Connection to Prior Work}
\label{sec:prior_work}

This framework builds directly on three foundational papers by the author:

\begin{enumerate}[leftmargin=2.2em]
  \item \textbf{``A Mechanical Route to $G$''} \cite{SkavbergMechG2025}: Derives the gravitational constant $G$ from the bulk modulus, mass density, and coherence length of a condensate substrate, establishing the mechanical-geometric linkage $G \sim \ell_{\mathrm{coh}}^2 / (K \rho)$.
  \item \textbf{``A Mechanical Extension of General Relativity: The Singularity in Equilibrium''} \cite{SkavbergMechGR2025}: Shows that a condensate in equilibrium reproduces Einstein--Hilbert dynamics, with the Big Bang singularity replaced by a nucleation-and-growth phase transition. Validates the framework against classical GR tests (perihelion precession, light deflection, Shapiro delay, frame dragging).
  \item \textbf{``Planck Field --- Resolving Hulse--Taylor Binary''} \cite{SkavbergPlanckField2025}: Applies the condensate framework to binary pulsar systems, demonstrating agreement with observed orbital decay rates and gravitational-wave emission timescales.
\end{enumerate}

These papers collectively establish the \emph{microphysical foundation} for gravitational response and validate the equilibrium limit against precision tests. The present manuscript extends this foundation to \emph{cosmological scales}, introducing boundary-layer kinematics and early-Universe phenomenology, and specifying rigorous validation criteria against CMB, BBN, BAO, SNe, and LSS data.

\TODO{PHYS-AGENT: Add references to analogue gravity reviews (Barceló et al., Volovik, etc.) and condensed-matter cosmology literature (e.g., Volovik's ``The Universe in a Helium Droplet'').}

\subsection{Falsifiability and Validation Checkpoints}
\label{sec:falsifiability}

A key strength of this framework is its \emph{falsifiability}. By design, the condensate model \emph{must} recover standard physics in all well-tested regimes:
\begin{itemize}[leftmargin=1.8em]
  \item \textbf{Classical GR tests} (Solar System, binary pulsars): The equilibrium condensate reduces to Einstein--Hilbert dynamics, preserving all PPN-parameter constraints.
  \item \textbf{Gravitational-wave propagation}: Tensor modes must propagate luminally at late times, satisfying the GW170817/GRB 170817A bound $|c_T/c - 1| \lesssim 10^{-15}$ \cite{GW170817}.
  \item \textbf{Big-Bang Nucleosynthesis (BBN)}: Any early-time modification to $H(z)$ must not alter the expansion rate at $T \sim \text{MeV}$ beyond the few-percent level constrained by light-element abundances \cite{BBN_Gbounds}.
  \item \textbf{CMB/Planck 2018}: The framework must reproduce Planck TT/TE/EE+lensing likelihoods at $\Lambda$CDM-like $\chi^2$, with acoustic peak phases and the damping tail preserved \cite{Planck2018}.
  \item \textbf{BAO and SNe}: Late-time observables (BAO sound horizon, SNe distance moduli) must remain internally consistent with CMB-inferred parameters.
  \item \textbf{Distance ladders}: The preferred parameter region must be compatible with local $H_0$ measurements from Cepheids, TRGB, JWST cross-calibrators, etc.\ \cite{RiessReview2024}.
  \item \textbf{Large-scale structure}: Growth history and weak lensing ($S_8$) should not be worsened relative to $\Lambda$CDM.
\end{itemize}

These ``validation checkpoints'' are formalized in Section \ref{sec:validation} with quantitative pass/fail criteria. Deviations from standard physics are permitted \emph{only} in a narrow early-time window (post-BBN, pre-recombination), where boundary-layer energy reservoirs can transiently raise $H(a)$ in an EDE-like manner \cite{EDEReview2023,PoulinReview2023}.

\GRAPHAGENT~\textbf{Figure 1 placeholder:} Schematic diagram of the two-phase system (protofluid exterior, condensate interior, conversion boundary $\Sigma$) with annotations for impedance mismatch, reflection/transmission coefficients, and the coarse-grained FLRW interior.

\subsection{Roadmap of This Paper}
\label{sec:roadmap_intro}

The remainder of this paper is organized as follows:
\begin{itemize}[leftmargin=1.8em]
  \item \textbf{Section \ref{sec:postulates}}: Core postulates defining the two-phase medium, phase-transition Big Bang, mechanical-geometric linkage, impedance mismatch, and causality/locality.
  \item \textbf{Section \ref{sec:kinematics}}: Kinematics of nucleation and growth, including Rankine--Hugoniot jump conditions at the conversion boundary, impedance-based reflection/transmission, and the coarse-grained FLRW expansion with a transient $\umech(a)$ bump.
  \item \textbf{Section \ref{sec:validation}}: Validation of foundational proofs---quantitative pass/fail criteria for GR tests, GW propagation, BBN, CMB/Planck, BAO/SNe, distance ladders, and LSS.
  \item \textbf{Section \ref{sec:relation}}: Relation to established results and discussion of why falsifiability is clear.
  \item \textbf{Section \ref{sec:roadmap}}: Research roadmap enumerating seven sequenced follow-on papers to stress-test the framework.
  \item \textbf{Section \ref{sec:conclusions}}: Conclusions and outlook.
\end{itemize}

\TODO{PHYS-AGENT: Verify that all cited references are consistent with the latest literature (2024--2025). Add missing citations for analogue gravity, EDE models, and $S_8$ tension.}

% ============== SECTION 2: CORE POSTULATES ==============
\section{Core Postulates}
\label{sec:postulates}

We formalize the foundational assumptions of the condensate cosmology framework as five core postulates. These postulates define the physical setting, the origin of the Big Bang, the mechanism for gravitational response, the role of boundary-layer physics, and the requirement of causality/locality.

\begin{postulate}[Two-Phase Medium]
\label{post:two_phase}
The Universe is embedded in a \emph{universal protofluid}---a pre-existing medium characterized by an effective equation of state, sound speed $\cs^{\mathrm{pf}}$, and impedance $\Zpf = \rho_{\mathrm{pf}} \cs^{\mathrm{pf}}$. At a critical event (the Big Bang), a \emph{coherent quantum condensate} nucleates within the protofluid and subsequently expands. The observable Universe is identified with the interior of this condensate. The two phases are separated by a codimension-one hypersurface $\Sigma$ (the \emph{conversion boundary}) of finite thickness, surface stress-energy, and non-equilibrium microphysics.
\end{postulate}

\begin{postulate}[Phase-Transition Origin of the Big Bang]
\label{post:phase_transition}
The Big Bang corresponds to the \emph{nucleation of a critical seed} at cosmic time $t = 0$, initiating a first-order-like phase transition. The conversion boundary $\Sigma$ propagates outward with proper speed $\uf(t)$, converting protofluid into condensate. The singularity of classical General Relativity is replaced by this nucleation event; the ``initial conditions'' of the Universe are the thermodynamic and quantum properties of the protofluid at the onset of nucleation.
\end{postulate}

\TODO{PHYS-AGENT: Expand discussion of nucleation theory (classical nucleation, quantum tunneling, bubble nucleation in field theory). Reference Coleman--De Luccia, Linde, etc.}

\begin{postulate}[Mechanical--Geometric Linkage]
\label{post:mech_geo}
The condensate's effective stiffness---characterized by bulk modulus $K$, shear modulus (if any), coherence length $\ell_{\mathrm{coh}}$, and mass density $\rho$---determines the gravitational constant $G$ and the form of the gravitational action. In equilibrium (i.e., far from the boundary, or after boundary effects have decayed), the effective action reduces to the Einstein--Hilbert form:
\begin{equation}
\label{eq:EH_action}
S_{\mathrm{eff}} = \frac{1}{16\pi G} \int \dx^4 x \sqrt{-g} \, R + S_{\mathrm{matter}}.
\end{equation}
The constitutive law connecting mechanical stiffness to geometric response is derived in \cite{SkavbergMechG2025,SkavbergMechGR2025}. Deviations from \eqref{eq:EH_action} can occur only during a short, early boundary epoch when non-equilibrium boundary-layer dynamics dominate.
\end{postulate}

\MATHAGENT~\textbf{TODO:} Provide a self-contained derivation of Equation \eqref{eq:EH_action} from the condensate constitutive law, or a clear reference to the relevant sections of \cite{SkavbergMechG2025,SkavbergMechGR2025}.

\begin{postulate}[Impedance Mismatch at the Boundary]
\label{post:impedance}
Linear perturbations (acoustic waves, gravitational waves, density fluctuations) propagating in the protofluid and condensate phases experience different effective impedances, $\Zpf$ and $\Zcond$ respectively. At the conversion boundary $\Sigma$, impedance mismatch leads to partial reflection and transmission, with coefficients
\begin{equation}
\label{eq:reflection_transmission}
R = \left| \frac{\Zcond - \Zpf}{\Zcond + \Zpf} \right|^2, \qquad T = \frac{4 \Zcond \Zpf}{(\Zcond + \Zpf)^2}, \qquad R + T = 1.
\end{equation}
A strong mismatch ($\Zcond \gg \Zpf$ or $\Zcond \ll \Zpf$) results in significant reflection, creating \emph{boundary-layer energy reservoirs} that transiently accumulate energy and can briefly raise the expansion rate $H(a)$. The amplitude and duration of this effect must be small and localized in time to satisfy observational constraints (see Section \ref{sec:validation}).
\end{postulate}

\MATHAGENT~\textbf{TODO:} Derive Equation \eqref{eq:reflection_transmission} from first principles (conservation of energy and momentum at the boundary, matching conditions for wave amplitudes). Include a discussion of Doppler-like shifts in the boundary frame and mode selection.

\begin{postulate}[Causality and Locality]
\label{post:causality}
Signal propagation in the condensate interior respects local light cones defined by the effective metric $g_{\mu\nu}$. Gravitational waves (tensor perturbations) propagate luminally at late times, $c_T = c$, satisfying the stringent GW170817/GRB 170817A bound $|c_T/c - 1| \lesssim 10^{-15}$ \cite{GW170817}. No superluminal signaling or acausal effects are permitted. During the early boundary epoch, transient modifications to $H(a)$ are local in time (confined to a narrow window in $\ln a$) and do not violate causality.
\end{postulate}

\TODO{PHYS-AGENT: Clarify the distinction between causal structure in the condensate interior versus across the boundary. Discuss whether the boundary propagation speed $\uf$ can exceed the local sound speed, and under what conditions this is permitted without violating causality.}

\GRAPHAGENT~\textbf{Figure 2 placeholder:} Visual representation of the five postulates: (a) two-phase schematic, (b) nucleation/growth timeline, (c) mechanical-to-geometric linkage diagram, (d) impedance mismatch and R/T coefficients, (e) causal light-cone structure.

% ============== SECTION 3: KINEMATICS OF NUCLEATION AND GROWTH ==============
\section{Kinematics of Nucleation and Growth}
\label{sec:kinematics}

This section develops the kinematic framework governing the conversion boundary and the resulting cosmological expansion. We first derive the generalized Rankine--Hugoniot jump conditions at the boundary, then analyze impedance-based reflection and transmission, and finally introduce the coarse-grained FLRW description with a transient $\umech(a)$ energy component.

\subsection{Conversion Boundary and Rankine--Hugoniot Relations}
\label{sec:RH}

Let $\Sigma$ denote the worldvolume of the conversion boundary, with unit spacelike normal $n_\mu$ (pointing from condensate to protofluid) and timelike 4-velocity $u_\mu^{\Sigma}$ satisfying $u_\mu^{\Sigma} u^{\mu \Sigma} = -1$ and $u_\mu^{\Sigma} n^\mu = 0$. The proper speed of the boundary in the protofluid frame is $\uf$. Define the jump of a bulk quantity $X$ across $\Sigma$ as
\begin{equation}
\label{eq:jump_def}
[X] \equiv X_{\mathrm{c}} - X_{\mathrm{pf}},
\end{equation}
where subscripts ``c'' and ``pf'' denote condensate and protofluid values, respectively.

Integral conservation of stress-energy across $\Sigma$ yields the generalized Rankine--Hugoniot relations:
\begin{equation}
\label{eq:RH_general}
\big[ T^{\mu\nu} n_\mu \big] = - \nabla_\alpha S^{\alpha\nu},
\end{equation}
where $T^{\mu\nu}$ is the bulk stress-energy tensor and $S^{\mu\nu}$ is the surface stress-energy tensor of the boundary layer. If $S^{\mu\nu}$ is conserved on $\Sigma$ (i.e., $\nabla_\alpha S^{\alpha\nu} = 0$ intrinsically), then the jump is zero:
\begin{equation}
\label{eq:RH_simple}
\big[ T^{\mu\nu} n_\mu \big] = 0.
\end{equation}

For a perfect fluid with $T^{\mu\nu} = (\rho + p) u^\mu u^\nu + p g^{\mu\nu}$, the normal component yields the standard shock-jump conditions relating density, pressure, and velocity discontinuities. In the condensate framework, the boundary layer can accumulate energy (non-zero $S^{\mu\nu}$), leading to a modified Hugoniot relation that sources transient energy reservoirs.

\MATHAGENT~\textbf{TODO:} Expand this subsection with explicit calculations:
\begin{itemize}
  \item Derive \eqref{eq:RH_general} from the Gauss--Stokes theorem and stress-energy conservation $\nabla_\mu T^{\mu\nu} = 0$ in each bulk phase.
  \item Write out the components of $S^{\mu\nu}$ in the boundary rest frame (surface tension, surface pressure, kinetic energy density).
  \item Show the connection to classical shock physics (Rankine--Hugoniot for gas dynamics, Taub adiabat, etc.).
  \item Analyze the limit where $S^{\mu\nu}$ is small (weak boundary layer) versus large (strong boundary layer).
\end{itemize}

\subsection{Impedance Mismatch, Reflection, and Transmission}
\label{sec:impedance_detail}

Linearized perturbations propagating in the protofluid and condensate phases are characterized by effective impedances:
\begin{equation}
\label{eq:impedances}
\Zpf = \rho_{\mathrm{pf}} \cs^{\mathrm{pf}}, \qquad \Zcond = \rho_{\mathrm{c}} \cs^{\mathrm{c}},
\end{equation}
where $\cs^{\mathrm{pf/c}}$ are the respective sound speeds. At the boundary $\Sigma$, impedance mismatch determines the reflection and transmission coefficients (as stated in Postulate \ref{post:impedance}):
\begin{equation}
\label{eq:RT_repeat}
R = \left| \frac{\Zcond - \Zpf}{\Zcond + \Zpf} \right|^2, \qquad T = \frac{4 \Zcond \Zpf}{(\Zcond + \Zpf)^2}.
\end{equation}

If $\Zcond \gg \Zpf$ (stiff condensate, soft protofluid), most energy is reflected back into the protofluid, with a small transmitted component entering the condensate. This reflection accumulates energy in the boundary layer, creating a transient energy reservoir. Conversely, if $\Zcond \ll \Zpf$, most energy is transmitted, and the boundary effect is minimal.

\paragraph{Doppler-like shifts and mode selection.}
In the boundary rest frame, incoming waves from the protofluid experience a Doppler shift due to the boundary's motion at speed $\uf$. Modes with frequency $\omega_{\mathrm{pf}}$ in the protofluid frame are shifted to
\begin{equation}
\label{eq:doppler_shift}
\omega_{\Sigma} = \omega_{\mathrm{pf}} \left(1 - \frac{\uf}{\cs^{\mathrm{pf}}} \right)^{-1}
\end{equation}
in the boundary frame (for head-on incidence). If $\uf \to \cs^{\mathrm{pf}}$, the boundary becomes supersonic, and a shock front forms. The reflected and transmitted waves then have frequencies determined by the impedance ratio and the boundary velocity, leading to frequency-dependent reflection/transmission and potential resonances that enhance energy accumulation at specific scales.

\MATHAGENT~\textbf{TODO:} Provide detailed derivations:
\begin{itemize}
  \item Matching conditions for wave amplitudes, pressures, and velocities at the boundary.
  \item Doppler-shift formula \eqref{eq:doppler_shift} from Lorentz transformations or acoustic analogue.
  \item Analysis of resonance conditions: at what $\omega_{\mathrm{pf}}$ and $\uf$ do we get maximum reflection and energy buildup?
  \item Connection to acoustic impedance in classical acoustics, shock-tube experiments, and condensed-matter interfaces.
\end{itemize}

\GRAPHAGENT~\textbf{Figure 3 placeholder:} Plots of $R(\Zcond/\Zpf)$ and $T(\Zcond/\Zpf)$ showing the transition from weak mismatch (small $R$) to strong mismatch (large $R$). Annotate regimes relevant to the condensate cosmology framework.

\subsection{Coarse-Grained FLRW Expansion with Transient Energy Component}
\label{sec:FLRW_coarse}

On scales much larger than the boundary thickness and coherence length, the condensate interior admits an effective Friedmann--Lema\^itre--Robertson--Walker (FLRW) description with scale factor $a(t)$, where $t$ is cosmic time. The Friedmann equation governing the expansion rate $H \equiv \dot{a}/a$ is
\begin{equation}
\label{eq:Friedmann_general}
H^2(a) = \frac{8\pi G_0}{3} \left[ \rhostd(a) + \umech(a) \right],
\end{equation}
where $\rhostd(a)$ is the standard energy density (radiation, matter, $\Lambda$) and $\umech(a)$ is a transient energy component sourced by boundary-layer dynamics.

\paragraph{Standard density evolution.}
The standard components evolve as:
\begin{equation}
\label{eq:rho_std}
\rhostd(a) = \rho_{\mathrm{r},0} a^{-4} + \rho_{\mathrm{m},0} a^{-3} + \rho_{\Lambda,0},
\end{equation}
where subscripts $r$, $m$, $\Lambda$ denote radiation, matter, and dark energy, and the subscript $0$ denotes present-day values.

\paragraph{Transient boundary-layer energy.}
The boundary-layer energy reservoir $\umech(a)$ is modeled phenomenologically as a Gaussian bump in $\ln a$:
\begin{equation}
\label{eq:u_mech}
\umech(a) = \rho_{\mathrm{c},0} \, \fmech \, \exp\!\left( -\frac{\ln^2(a/\ac)}{2\sigmamech^2} \right),
\end{equation}
where:
\begin{itemize}[leftmargin=1.8em]
  \item $\rho_{\mathrm{c},0} = 3H_0^2/(8\pi G_0)$ is the critical density today,
  \item $\fmech \ll 1$ is the peak fractional energy density (peak value of $\umech/\rhocrit$ at $a = \ac$),
  \item $\ac$ is the scale factor at which the bump peaks (typically early times, $\ac \sim 10^{-3}$--$10^{-2}$),
  \item $\sigmamech$ is the width of the bump in $\ln a$ (narrow, $\sigmamech \sim 0.1$--$0.3$).
\end{itemize}

This parameterization is analogous to Early Dark Energy (EDE) models \cite{EDEReview2023,PoulinReview2023}, but is \emph{derived} from boundary-layer impedance physics rather than postulated as an additional scalar field. The parameters $(\fmech, \ac, \sigmamech)$ are related to the impedance ratio $\Zcond/\Zpf$, boundary speed $\uf$, and boundary thickness via a mapping to be derived in a follow-on paper (see Section \ref{sec:roadmap}, Paper 4).

\MATHAGENT~\textbf{TODO:} Derive the functional form \eqref{eq:u_mech} from:
\begin{itemize}
  \item Conservation of energy at the boundary: $\partial_t (\rho_{\mathrm{boundary}} V_{\mathrm{boundary}}) + \dots$
  \item Impedance-based energy accumulation rate: $\dot{E}_{\mathrm{acc}} \propto R \times (\text{incident flux})$.
  \item Decay timescale for boundary-layer energy (dissipation, dilution due to expansion).
  \item Dimensional analysis and scaling arguments to fix $\fmech$, $\ac$, $\sigmamech$ in terms of $(\Zpf, \Zcond, \uf, \ell_{\mathrm{boundary}})$.
\end{itemize}

\paragraph{Equation of state of $\umech$.}
The effective equation-of-state parameter $w_{\mathrm{mech}} \equiv p_{\mathrm{mech}}/\umech$ can be inferred from the time evolution of $\umech$. Energy conservation $\dot{\rho} + 3H(\rho + p) = 0$ implies
\begin{equation}
\label{eq:w_mech}
w_{\mathrm{mech}} = -1 - \frac{1}{3H} \frac{\dx \ln \umech}{\dx t}.
\end{equation}
For the Gaussian form \eqref{eq:u_mech}, $w_{\mathrm{mech}}$ transitions from $\approx -1/3$ (radiation-like) at early times to more negative values near the peak, then back to $\approx -1/3$ at late times (as the bump decays away). This time-varying equation of state is a key signature distinguishing boundary-layer energy from a cosmological constant.

\MATHAGENT~\textbf{TODO:} Calculate $w_{\mathrm{mech}}(a)$ explicitly from \eqref{eq:u_mech} and \eqref{eq:w_mech}. Plot $w_{\mathrm{mech}}$ versus $a$ for representative parameter values. Discuss implications for sound speed $\cs^2 = \dx p_{\mathrm{mech}}/\dx \umech$ and perturbation growth.

\GRAPHAGENT~\textbf{Figure 4 placeholder:} Plot of $\umech(a)/\rhocrit(a)$ versus $a$ (or redshift $z$) for representative $(\fmech, \ac, \sigmamech)$. Overlay standard $\Lambda$CDM density components (radiation, matter, $\Lambda$) to show the relative importance of the boundary-layer energy bump.

\GRAPHAGENT~\textbf{Figure 5 placeholder:} Plot of $H(a)/H_{\Lambda\mathrm{CDM}}(a)$ versus $a$ showing the percent-level enhancement during the boundary epoch.

% ============== SECTION 4: VALIDATION OF FOUNDATIONAL PROOFS ==============
\section{Validation of Foundational Proofs (Pass/Fail Guardrails)}
\label{sec:validation}

This section establishes quantitative pass/fail criteria for the condensate framework against well-established observational pillars. The framework \emph{must} satisfy these guardrails to be considered viable; failure in any single test falsifies the model. Each subsection specifies:
\begin{itemize}[leftmargin=1.8em]
  \item The \emph{limit} in which standard physics must be recovered,
  \item The \emph{quantitative criterion} constituting a pass,
  \item The \emph{rationale} for why this test is non-negotiable.
\end{itemize}

\subsection{Classical GR Tests (Solar System, Binary Pulsars)}
\label{sec:GR_tests}

\paragraph{Test statement.}
In the equilibrium condensate (no boundary effects, $\umech \to 0$), the effective action must reduce to the Einstein--Hilbert form \eqref{eq:EH_action}, recovering all classical General Relativity predictions: perihelion precession of Mercury, light deflection by the Sun, Shapiro time delay, frame dragging (Lense--Thirring effect), and binary pulsar orbital decay (Hulse--Taylor PSR B1913+16, Double Pulsar PSR J0737-3039).

\paragraph{Pass criterion.}
No deviations from Parameterized Post-Newtonian (PPN) constraints at the quoted experimental precision. Specifically:
\begin{itemize}[leftmargin=1.8em]
  \item PPN parameters: $|\gamma_{\mathrm{PPN}} - 1| < 2.3 \times 10^{-5}$, $|\beta_{\mathrm{PPN}} - 1| < 8 \times 10^{-5}$ (Cassini spacecraft) \cite{PPN_Cassini}.
  \item Mercury perihelion: prediction matches observation to $\sim 0.01\%$ \cite{PPN_Perihelion}.
  \item Gravitational time delay: agreement with Cassini/Viking to $\sim 10^{-5}$ \cite{PPN_Shapiro}.
  \item Binary pulsar decay: Hulse--Taylor orbital decay matches GR prediction to $\sim 0.2\%$ \cite{HulseTaylor}; Double Pulsar to $\sim 0.05\%$ \cite{DoublePulsar}.
\end{itemize}

\paragraph{Rationale.}
These tests probe the gravitational field in the weak-field, slow-motion limit and in strong-field, relativistic regimes (neutron stars). Any residual modification from boundary-layer physics would be cosmologically red-shifted away by the time of structure formation; these tests occur in the present-day equilibrium condensate where $\umech = 0$.

\TODO{PHYS-AGENT: Verify that the cited PPN bounds are up-to-date (2024--2025). Add references to recent Solar System ephemeris tests (e.g., BepiColombo, Juno gravity experiments).}

\subsection{Gravitational-Wave Propagation}
\label{sec:GW_prop}

\paragraph{Test statement.}
The tensor polarization modes of gravitational waves (GWs) must propagate luminally in the present-day Universe, respecting the stringent bound from the binary neutron star merger GW170817 and its electromagnetic counterpart GRB 170817A \cite{GW170817}.

\paragraph{Pass criterion.}
The speed of tensor modes $c_T$ must satisfy
\begin{equation}
\label{eq:cT_bound}
\left| \frac{c_T}{c} - 1 \right| \lesssim 10^{-15}
\end{equation}
at late times ($z \lesssim 0.01$). No anomalous dispersion, extra polarization states, or frequency-dependent propagation speeds are permitted.

\paragraph{Rationale.}
The GW170817 bound is one of the tightest constraints on modified gravity theories. The condensate framework naturally satisfies this bound because:
\begin{itemize}[leftmargin=1.8em]
  \item The boundary-layer energy $\umech(a)$ is localized at early times ($a \sim \ac \sim 10^{-3}$--$10^{-2}$) and has decayed to zero by late times.
  \item In equilibrium, tensor modes propagate on the effective metric $g_{\mu\nu}$ with no additional damping or dispersion.
  \item The Planck-field mechanism \cite{SkavbergPlanckField2025} ensures that GW emission from compact binaries matches GR predictions.
\end{itemize}

\TODO{PHYS-AGENT: Expand discussion of modified dispersion relations (e.g., massive graviton, Lorentz violation) and why the condensate framework avoids them. Reference recent LIGO/Virgo/KAGRA constraints on $c_T$ and polarization modes.}

\subsection{Big-Bang Nucleosynthesis (BBN)}
\label{sec:BBN}

\paragraph{Test statement.}
The primordial abundances of light elements (deuterium D, helium-3 $^3$He, helium-4 $^4$He, lithium-7 $^7$Li) are sensitive to the expansion rate $H(T)$ at temperatures $T \sim 0.01$--$1$ MeV, corresponding to times $t \sim 1$--$200$ seconds after the Big Bang. Any modification to $H(T)$ in this epoch directly affects the neutron-to-proton freeze-out ratio and the subsequent nuclear reaction rates.

\paragraph{Pass criterion.}
The best-fit parameters $(\fmech, \ac, \sigmamech)$ that address the $H_0$ tension must be consistent with BBN-derived bounds on the effective gravitational constant $G_{\mathrm{BBN}}$ and expansion rate $H_{\mathrm{BBN}}$:
\begin{equation}
\label{eq:BBN_G_bound}
\left| \frac{G_{\mathrm{BBN}}}{G_0} - 1 \right| \lesssim 0.05 \quad \text{(95\% CL)},
\end{equation}
as derived from joint fits to D/H, $^4$He, and $^7$Li \cite{BBN_Gbounds,PDG_BBN}. Equivalently, the expansion rate during BBN must not deviate by more than $\sim 5\%$ from the standard $\Lambda$CDM prediction.

\paragraph{Rationale.}
BBN is one of the earliest and most robust observational probes of cosmology. The measured deuterium abundance D/H from quasar absorption systems is precise to $\sim 1\%$ \cite{Cooke_D2014}, and the $^4$He mass fraction $Y_p$ is measured to $\sim 1\%$ from low-metallicity HII regions \cite{AdesOlivares_He2020}. These measurements constrain the baryon-to-photon ratio $\eta = n_b/n_\gamma$ and the expansion rate $H(T_{\mathrm{BBN}})$ very tightly. Any early-time boundary-layer energy must:
\begin{itemize}[leftmargin=1.8em]
  \item Be negligible at $a \sim 10^{-9}$--$10^{-8}$ (corresponding to $T \sim 0.1$--$1$ MeV),
  \item OR be centered at much earlier times ($\ac \ll 10^{-8}$) and have decayed away by BBN,
  \item OR be centered at later times ($\ac \gg 10^{-8}$) and not yet turned on during BBN.
\end{itemize}

The narrow Gaussian form \eqref{eq:u_mech} with peak at $\ac \sim 10^{-3}$--$10^{-2}$ (corresponding to $z \sim 10^2$--$10^3$, i.e., post-recombination but pre-structure-formation) naturally satisfies this requirement, as the bump is far in the future relative to BBN.

\MATHAGENT~\textbf{TODO:} Calculate $\umech(a_{\mathrm{BBN}})/\rhocrit(a_{\mathrm{BBN}})$ for representative $(\fmech, \ac, \sigmamech)$ and show it is $\ll 0.01$ (i.e., negligible at BBN). Integrate the modified Friedmann equation to compute $H(a_{\mathrm{BBN}})$ and verify the bound \eqref{eq:BBN_G_bound} is satisfied.

\TODO{PHYS-AGENT: Update BBN references to latest Planck/PDG/CMB-S4 compilations. Discuss lithium problem and whether boundary-layer physics offers any resolution (likely not, as lithium abundance is mainly a nuclear physics issue).}

\subsection{CMB/Planck 2018 + BAO + SNe (Acoustic Proofs)}
\label{sec:CMB_Planck}

\paragraph{Test statement.}
The Cosmic Microwave Background (CMB) power spectra (temperature TT, polarization EE/TE, lensing) from Planck 2018 \cite{Planck2018}, combined with Baryon Acoustic Oscillations (BAO) from large-scale galaxy surveys \cite{BOSS_BAO,eBOSS_BAO} and Type Ia supernovae (SNe) distance moduli \cite{Pantheon,PantheonPlus}, constitute the ``acoustic proofs'' of standard cosmology. These observables are extremely sensitive to the expansion history $H(z)$, the sound horizon at decoupling $\rs$, and the angular diameter distance $d_A(z)$.

\paragraph{Pass criterion.}
The condensate framework must:
\begin{enumerate}[leftmargin=2.2em]
  \item Achieve a $\chi^2$ fit to Planck 2018 TT/TE/EE+lensing likelihoods that is competitive with $\Lambda$CDM: $\Delta\chi^2_{\mathrm{CMB}} \lesssim 5$ (or $\chi^2/\mathrm{dof} \sim 1$).
  \item Preserve the CMB acoustic peak phases (multipole positions $\ell_1, \ell_2, \dots$) and damping tail to within Planck statistical uncertainties.
  \item Remain consistent with BAO measurements of $d_A(z)/\rs$ and $H(z) \rs$ at $z \sim 0.1$--$2.3$ within $1\sigma$.
  \item Remain consistent with SNe distance moduli $\mu(z)$ within $1\sigma$.
  \item If the framework addresses the $H_0$ tension by reducing $\rs$ via the boundary-layer bump, the reduction must be modest ($\sim 2$--$5\%$) and not degrade BAO/SNe fits.
\end{enumerate}

\paragraph{Rationale.}
The CMB acoustic structure is the most precise ``standard ruler'' in cosmology, with the sound horizon at decoupling measured to $\sim 0.2\%$ \cite{Planck2018}. The angular positions of the acoustic peaks encode the geometry and expansion history from decoupling ($z \sim 1100$) to today. Any modification to $H(z)$ that alters $\rs$ will shift the peak positions; if the shift is too large, the CMB power spectrum will be inconsistent with Planck data.

The boundary-layer energy bump $\umech(a)$ is designed to:
\begin{itemize}[leftmargin=1.8em]
  \item Peak at $\ac \sim 10^{-3}$--$10^{-2}$ (i.e., $z \sim 10^2$--$10^3$, between recombination and today),
  \item Have small amplitude $\fmech \lesssim 0.1$ (i.e., $\lesssim 10\%$ of the critical density at peak),
  \item Have narrow width $\sigmamech \sim 0.1$--$0.3$ in $\ln a$ (brief duration in cosmic time),
\end{itemize}
such that the integral effect on $\rs$ is a few percent, raising the CMB-inferred $H_0$ closer to the distance-ladder value \cite{EDEReview2023,PoulinReview2023}, while leaving the peak phases and damping tail nearly unchanged.

\MATHAGENT~\textbf{TODO:} Implement the modified Friedmann equation \eqref{eq:Friedmann_general} with \eqref{eq:u_mech} in a CMB Boltzmann code (e.g., CLASS, CAMB). Run parameter scans over $(\fmech, \ac, \sigmamech)$ and compute:
\begin{itemize}
  \item The sound horizon at decoupling: $\rs = \int_0^{t_{\mathrm{dec}}} \cs(t) \, \dx t / a(t_{\mathrm{dec}})$.
  \item The CMB power spectra $C_\ell^{TT}$, $C_\ell^{TE}$, $C_\ell^{EE}$, and compare to Planck 2018 data.
  \item The best-fit $H_0$ and $\Omega_m$ from CMB+BAO+SNe joint fits.
  \item The $\chi^2$ improvement (or degradation) relative to $\Lambda$CDM.
\end{itemize}

\TODO{PHYS-AGENT: Update BAO and SNe references to latest releases (DESI 2024, DES-SN5YR, Pantheon+, etc.). Discuss whether the condensate framework can simultaneously address the $H_0$ and $S_8$ tensions, or whether they require separate mechanisms.}

\GRAPHAGENT~\textbf{Figure 6 placeholder:} CMB temperature power spectrum $C_\ell^{TT}$ for $\Lambda$CDM (blue) versus condensate framework with boundary-layer bump (red), showing acoustic peak positions and damping tail. Overlay Planck 2018 data points with error bars.

\GRAPHAGENT~\textbf{Figure 7 placeholder:} Joint constraints on $(H_0, \Omega_m)$ from CMB, BAO, and SNe, comparing $\Lambda$CDM (gray contours) and condensate framework (red contours). Show distance-ladder $H_0$ constraints (SH0ES, TRGB, etc.) as horizontal bands.

\subsection{Distance Ladders and Cross-Checks}
\label{sec:distance_ladders}

\paragraph{Test statement.}
Local measurements of the Hubble constant $H_0$ from distance ladders---using Cepheid variables, Tip of the Red Giant Branch (TRGB), Mira variables, Surface Brightness Fluctuations (SBF), and recently JWST cross-calibrations \cite{RiessReview2024,FreedmanTRGB2020}---provide independent constraints on $H_0$ that are in tension with the CMB-inferred value. The condensate framework aims to resolve this tension by modifying the early expansion history (via $\umech$), raising the CMB-inferred $H_0$ without introducing new late-time physics.

\paragraph{Pass criterion.}
The preferred parameter region $(\fmech, \ac, \sigmamech)$ from CMB+BAO+SNe fits must yield a best-fit $H_0$ that is:
\begin{enumerate}[leftmargin=2.2em]
  \item Consistent with the latest distance-ladder compilations at $\lesssim 2\sigma$ (e.g., SH0ES: $H_0 = 73.04 \pm 1.04$ km/s/Mpc \cite{Riess2022}; TRGB: $H_0 = 69.8 \pm 1.7$ km/s/Mpc \cite{FreedmanTRGB2020}).
  \item Compatible with JWST cross-checks (e.g., Cepheids in NGC 4258, Cepheids+TRGB in the LMC) \cite{JWST_Cepheids2023}.
  \item Does not reintroduce late-time anomalies (e.g., modifications to $H(z)$ at $z < 1$ that would be inconsistent with BAO or SNe).
\end{enumerate}

\paragraph{Rationale.}
The $H_0$ tension is the primary motivation for exploring early-time modifications to $\Lambda$CDM. If the condensate framework's boundary-layer bump successfully raises $H_0$ to $\sim 71$--$73$ km/s/Mpc while preserving CMB+BAO+SNe consistency, this would constitute strong evidence in favor of the framework. Conversely, if the bump degrades fits or introduces new tensions, the framework is falsified.

\TODO{PHYS-AGENT: Update distance-ladder references to 2024--2025 releases (e.g., SH0ES with JWST, Chicago pipeline TRGB updates, CCHP Cepheid recalibrations). Discuss systematic uncertainties (metallicity dependence, crowding, extinction).}

\subsection{Large-Scale Structure and Weak Lensing}
\label{sec:LSS}

\paragraph{Test statement.}
The growth of structure from initial density perturbations to the present-day matter power spectrum $P(k)$ and weak lensing shear power spectrum $C_\ell^{\kappa\kappa}$ is governed by the linearized Einstein equations (or equivalently, the growth equation in the Newtonian gauge). The amplitude of clustering, parameterized by $S_8 \equiv \sigma_8 (\Omega_m/0.3)^{0.5}$, is currently in tension between Planck CMB ($S_8 \sim 0.83$) and weak lensing surveys (KiDS, DES, HSC: $S_8 \sim 0.76$--$0.79$) \cite{PlanckS8,DES2021,KiDS2020}.

\paragraph{Pass criterion.}
The condensate framework must:
\begin{enumerate}[leftmargin=2.2em]
  \item Not worsen the $S_8$ tension: the best-fit $S_8$ from CMB+LSS should remain within the current $\sim 2\sigma$ discrepancy band.
  \item Preserve the shape of the matter power spectrum $P(k)$ on BAO scales ($k \sim 0.01$--$0.2$ $h$/Mpc) to $\lesssim 5\%$.
  \item Not produce unacceptable excess small-scale power ($k > 1$ $h$/Mpc) that would conflict with Lyman-$\alpha$ forest constraints \cite{LymanAlpha}.
  \item Remain consistent with CMB lensing power spectrum from Planck and ACT/SPT \cite{PlanckLensing,ACTLensing}.
\end{enumerate}

\paragraph{Rationale.}
The boundary-layer energy bump $\umech(a)$ can affect structure growth in two ways:
\begin{itemize}[leftmargin=1.8em]
  \item \textbf{Early-time effect:} If the bump peaks at $\ac \sim 10^{-3}$--$10^{-2}$ (i.e., $z \sim 10^2$--$10^3$), it occurs during matter domination but before most small-scale structure has collapsed. The enhanced $H(a)$ slightly suppresses growth during the bump epoch, which could in principle reduce $\sigma_8$ and alleviate the $S_8$ tension.
  \item \textbf{Initial conditions:} If the bump is placed even earlier (e.g., near matter-radiation equality), it could imprint features in the initial power spectrum $P_{\mathrm{init}}(k)$ by altering the sound horizon and transfer function.
\end{itemize}

However, most EDE models \emph{increase} $S_8$ slightly (by raising $\Omega_m$), making the tension worse \cite{HillEDE2020}. The condensate framework must be carefully tuned to avoid this pitfall.

\MATHAGENT~\textbf{TODO:} Run linear perturbation codes (CLASS/CAMB) with the modified background $H(a)$ from \eqref{eq:Friedmann_general} and compute:
\begin{itemize}
  \item The matter power spectrum $P(k)$ at $z = 0$.
  \item The growth factor $D(a)$ and growth rate $f(a) = \dx \ln D / \dx \ln a$.
  \item The weak lensing shear power spectrum $C_\ell^{\kappa\kappa}$.
  \item The derived $\sigma_8$ and $S_8$ from best-fit cosmologies.
\end{itemize}

\TODO{PHYS-AGENT: Update LSS references to DESI 2024, Euclid early results, Vera Rubin LSST forecasts. Discuss whether the condensate framework can address the $S_8$ tension or whether this requires additional modifications (e.g., neutrino masses, interacting dark energy).}

\GRAPHAGENT~\textbf{Figure 8 placeholder:} Matter power spectrum $P(k)$ for $\Lambda$CDM versus condensate framework, showing BAO wiggles and small-scale damping. Overlay data from BOSS, eBOSS, DES.

\GRAPHAGENT~\textbf{Figure 9 placeholder:} Joint constraints on $(S_8, H_0)$ from CMB, weak lensing, and BAO, showing the $\Lambda$CDM tension and the condensate framework's preferred region.

% ============== SECTION 5: RELATION TO ESTABLISHED RESULTS ==============
\section{Relation to Established Results and Why Falsifiability is Clear}
\label{sec:relation}

This section synthesizes the validation checkpoints from Section \ref{sec:validation} and clarifies how the condensate framework achieves falsifiability by design.

\subsection{Einstein--Hilbert Limit}
\label{sec:EH_limit}

In equilibrium---defined as the regime where the conversion boundary has long since passed, the boundary-layer energy has dissipated ($\umech \to 0$), and the condensate has settled into a stationary state---the effective action reduces exactly to the Einstein--Hilbert form \eqref{eq:EH_action}. This is a consequence of Postulate \ref{post:mech_geo}, which links the condensate's mechanical stiffness to the gravitational response via the constitutive law derived in \cite{SkavbergMechG2025,SkavbergMechGR2025}.

As a result, all classical tests of General Relativity---perihelion precession, light deflection, Shapiro delay, frame dragging, binary pulsar decay---are automatically recovered. The framework does \emph{not} introduce new PPN parameters or modify the gravitational field equations in the weak-field, slow-motion limit. This is a non-negotiable requirement: any residual deviation from GR in the Solar System or binary pulsar regime would falsify the framework.

\TODO{PHYS-AGENT: Expand this subsection with a detailed comparison to alternative theories (f(R) gravity, scalar-tensor theories, DGP, Horndeski). Emphasize that the condensate framework is \emph{not} a modified gravity theory in the traditional sense, but rather an emergent-gravity framework that reduces to GR in equilibrium.}

\subsection{Planck ``Proofs''}
\label{sec:Planck_proofs}

The CMB acoustic structure is often referred to as a ``proof'' of standard cosmology, in the sense that the acoustic peak positions, damping tail, and lensing power spectrum are all precisely predicted by $\Lambda$CDM with a handful of parameters. Any modification to the expansion history $H(z)$ that alters the sound horizon $\rs$ or the angular diameter distance $d_A(z_*)$ (where $z_* \sim 1100$ is the redshift of decoupling) will shift the peak positions and potentially degrade the fit to Planck data.

The condensate framework's boundary-layer energy bump $\umech(a)$ is designed to be:
\begin{itemize}[leftmargin=1.8em]
  \item \textbf{Brief:} The width $\sigmamech \sim 0.1$--$0.3$ in $\ln a$ corresponds to a short duration in cosmic time, $\Delta t / t \sim \sigmamech \sim 10\%$--$30\%$ of the age of the Universe at the time of the bump.
  \item \textbf{Small:} The peak amplitude $\fmech \lesssim 0.1$ ensures that the fractional change in $H(a)$ is at most $\sim 5\%$--$10\%$ at the peak.
  \item \textbf{Localized:} The peak is at $\ac \sim 10^{-3}$--$10^{-2}$, i.e., $z \sim 10^2$--$10^3$, which is after recombination ($z \sim 1100$) but before structure formation ($z \sim 10$).
\end{itemize}

These design choices ensure that the bump's integral effect on $\rs$ is modest ($\sim 2$--$5\%$ reduction), raising the CMB-inferred $H_0$ by a few km/s/Mpc, while leaving the acoustic peak \emph{phases} (which depend on the ratio $\ell_{\mathrm{peak}} \propto d_A(z_*)/\rs$) nearly unchanged. The damping tail, which depends on the photon diffusion scale, is also preserved because the bump does not significantly alter the baryon-photon coupling or the Silk damping scale.

This is a stringent falsifiability criterion: if CMB+BAO+SNe fits with the condensate framework yield $\Delta\chi^2 > 10$ or introduce new peak-phase discrepancies, the framework is ruled out.

\TODO{PHYS-AGENT: Add discussion of ``derived parameters'' from Planck (e.g., $\tau$ optical depth, $A_s$ primordial amplitude, $n_s$ spectral index) and whether the condensate framework affects them. Reference recent ACT and SPT high-$\ell$ CMB analyses.}

\subsection{BBN \& GW Guardrails}
\label{sec:BBN_GW_guardrails}

The two sharpest guardrails for early-Universe modifications are:

\paragraph{(i) Big-Bang Nucleosynthesis.}
BBN occurs at $t \sim 1$--$200$ seconds, corresponding to $T \sim 0.01$--$1$ MeV and $a \sim 10^{-9}$--$10^{-8}$. At these early times, the boundary-layer energy bump is either (1) not yet activated (if $\ac \gg 10^{-8}$), or (2) negligibly small (if $\ac \sim 10^{-3}$ and $\sigmamech \sim 0.2$, the tail of the Gaussian at $a \sim 10^{-8}$ is $\exp(-\ln^2(10^5)/2 \times 0.2^2) \sim \exp(-1.4 \times 10^{11}) \approx 0$). Thus, $H(a_{\mathrm{BBN}}) = H_{\Lambda\mathrm{CDM}}(a_{\mathrm{BBN}})$ to extremely high precision, and the BBN bound \eqref{eq:BBN_G_bound} is automatically satisfied.

\paragraph{(ii) Gravitational-wave speed.}
The GW170817 bound $|c_T/c - 1| \lesssim 10^{-15}$ is the tightest constraint on tensor-mode propagation. In the condensate framework, tensor modes are perturbations of the effective metric $g_{\mu\nu}$ and propagate on null geodesics in equilibrium. The boundary-layer energy $\umech(a)$ is a scalar (isotropic) contribution to the stress-energy and does not couple to tensor modes at linear order. Moreover, $\umech(a)$ has decayed to zero by late times ($z \lesssim 1$), so present-day GW propagation is governed purely by the equilibrium Einstein--Hilbert dynamics. Thus, $c_T = c$ exactly at $z \sim 0.01$ (the redshift of GW170817), and the bound is satisfied.

These two guardrails together confine the allowed parameter space to a narrow window: the bump must be post-BBN ($\ac \gg 10^{-8}$) and pre-structure-formation ($\ac \lesssim 10^{-1}$), with decay complete by $z \sim 1$. This is a highly falsifiable prediction: any detection of tensor-mode anomalies or BBN inconsistencies would rule out the framework.

\TODO{PHYS-AGENT: Discuss potential loopholes or fine-tuning concerns. Could the boundary-layer energy have a small residual tail at late times that is undetectable by current GW experiments but might be probed by future space-based detectors (LISA, BBO)? What about primordial GWs from inflation---does the condensate framework predict any modifications to the primordial tensor power spectrum?}

\subsection{Summary: Why Falsifiability is Clear}
\label{sec:falsifiability_summary}

The condensate framework is falsifiable at multiple levels:
\begin{enumerate}[leftmargin=2.2em]
  \item \textbf{Classical GR tests:} Any deviation from PPN constraints or binary pulsar predictions falsifies the framework.
  \item \textbf{GW propagation:} Any detection of $c_T \neq c$ at late times falsifies the framework.
  \item \textbf{BBN:} Any $\sim 5\%$ deviation in $H(T_{\mathrm{BBN}})$ that conflicts with light-element abundances falsifies the framework.
  \item \textbf{CMB:} Any $\Delta\chi^2 > 10$ degradation in Planck fits, or any new peak-phase discrepancies, falsifies the framework.
  \item \textbf{BAO/SNe:} Any inconsistency with late-time distance measurements falsifies the framework.
  \item \textbf{LSS:} Any worsening of the $S_8$ tension beyond current levels falsifies the framework.
\end{enumerate}

Conversely, if the framework passes all these tests \emph{and} reduces the $H_0$ tension, this would constitute strong evidence in favor of the condensate paradigm.

% ============== SECTION 6: RESEARCH ROADMAP ==============
\section{Research Roadmap (Sequenced Follow-On Papers)}
\label{sec:roadmap}

To rigorously test the condensate framework, we propose a sequenced series of seven follow-on papers, each addressing a specific aspect of the theory and its empirical validation. These papers are designed to be modular and self-contained, allowing parallel progress by different members of the research team.

\subsection{Paper 1: Early-Universe Baseline Fits}
\label{sec:paper1}

\paragraph{Objective.}
Implement the transient energy component $\umech(a)$ from Equation \eqref{eq:u_mech} in a CMB Boltzmann code (CLASS or CAMB) and perform parameter scans over $(\fmech, \ac, \sigmamech)$ to find the best-fit cosmology that maximizes consistency with Planck 2018 TT/TE/EE+lensing, BAO, SNe, and distance-ladder $H_0$ measurements.

\paragraph{Key deliverables.}
\begin{itemize}[leftmargin=1.8em]
  \item Modified Friedmann equation solver integrated into CLASS/CAMB.
  \item MCMC chains exploring $\{H_0, \Omega_b, \Omega_c, A_s, n_s, \tau, \fmech, \ac, \sigmamech\}$ parameter space.
  \item Best-fit $H_0$ and comparison to SH0ES, TRGB, and CMB-only values.
  \item $\Delta\chi^2$ relative to $\Lambda$CDM for CMB, BAO, SNe individually and jointly.
  \item Constraints on $(\fmech, \ac, \sigmamech)$ from CMB+BAO+SNe likelihoods.
\end{itemize}

\paragraph{Pass criterion.}
Improvement in $H_0$ tension (reduction from $\sim 5\sigma$ to $\lesssim 3\sigma$) without degrading CMB+BAO+SNe fits ($\Delta\chi^2 \lesssim 5$).

\MATHAGENT~\textbf{TODO:} Set up the CLASS/CAMB modification, define priors on $(\fmech, \ac, \sigmamech)$, and run preliminary MCMC chains to assess feasibility.

\subsection{Paper 2: BBN and Time-Variation Checks}
\label{sec:paper2}

\paragraph{Objective.}
Demonstrate that the preferred parameter region from Paper 1 is fully consistent with Big-Bang Nucleosynthesis constraints on $G_{\mathrm{BBN}}$ and the expansion rate at $T \sim 0.1$--$1$ MeV. Verify that $\umech(a_{\mathrm{BBN}})$ is negligibly small and that no conflict arises with bounds on early-time $G$ variation \cite{BBN_Gbounds}.

\paragraph{Key deliverables.}
\begin{itemize}[leftmargin=1.8em]
  \item Calculation of $\umech(a)$ evaluated at $a_{\mathrm{BBN}} \sim 10^{-9}$--$10^{-8}$ for best-fit parameters.
  \item Integration of modified Friedmann equation to compute $H(T_{\mathrm{BBN}})$ and compare to standard BBN predictions.
  \item Joint fit to D/H, $^4$He, $^7$Li abundances using modified $H(T)$; derive constraints on $G_{\mathrm{BBN}}/G_0$.
  \item Verification that $|G_{\mathrm{BBN}}/G_0 - 1| \lesssim 0.05$ is satisfied.
\end{itemize}

\paragraph{Pass criterion.}
The condensate framework's best-fit parameters lie entirely within the BBN-allowed region at 95\% CL.

\MATHAGENT~\textbf{TODO:} Interface with BBN codes (e.g., PRIMAT, AlterBBN) to compute light-element yields with the modified expansion history.

\subsection{Paper 3: GW Propagation and Causality}
\label{sec:paper3}

\paragraph{Objective.}
Prove that tensor modes propagate luminally ($c_T = c$) at late times in the condensate framework, satisfying the GW170817 bound \eqref{eq:cT_bound}. Analyze whether any transient modifications to $c_T$ or dispersion relations occur during the boundary-layer epoch, and if so, whether they are observable in future GW detectors (LISA, Einstein Telescope, Cosmic Explorer).

\paragraph{Key deliverables.}
\begin{itemize}[leftmargin=1.8em]
  \item Derivation of the tensor-mode equation of motion in the presence of $\umech(a)$.
  \item Calculation of $c_T(a)$ versus scale factor; proof that $c_T(a \to 1) = c$ exactly.
  \item Analysis of frequency-dependent dispersion: $c_T(f)$ for GW frequencies $f \sim 10^{-4}$--$10^3$ Hz.
  \item Constraints on anomalous polarization states (vector, scalar) from Stokes parameters of GW waveforms.
  \item Forecasts for LISA sensitivity to residual $c_T$ deviations at $z \sim 1$--$10$.
\end{itemize}

\paragraph{Pass criterion.}
$c_T/c = 1$ to within $10^{-15}$ at $z \lesssim 1$, and no detectable dispersion or polarization anomalies in current or near-future GW data.

\MATHAGENT~\textbf{TODO:} Solve the linearized Einstein equations for tensor modes on the modified FLRW background. Compute the time-dependent effective mass or friction term induced by $\umech(a)$.

\TODO{PHYS-AGENT: Survey recent GW theory literature on modified dispersion, extra polarizations, and parity violation. Reference LIGO/Virgo/KAGRA parameterized tests of GR.}

\subsection{Paper 4: Boundary-Layer Kinematics and Wave Compression}
\label{sec:paper4}

\paragraph{Objective.}
Derive the boundary-layer energy reservoir $\umech(a)$ from first principles, starting from the impedance mismatch, reflection coefficient $R$, boundary speed $\uf$, and boundary thickness $\ell_{\mathrm{boundary}}$. Establish the mapping from microscopic boundary parameters $(\Zpf, \Zcond, \uf, \ell_{\mathrm{boundary}})$ to phenomenological parameters $(\fmech, \ac, \sigmamech)$.

\paragraph{Key deliverables.}
\begin{itemize}[leftmargin=1.8em]
  \item Rankine--Hugoniot analysis with surface stress-energy $S^{\mu\nu}$ included.
  \item Impedance-based energy accumulation rate: $\dot{E}_{\mathrm{acc}} = R \times \Phi_{\mathrm{incident}}$ (incident energy flux).
  \item Time evolution of boundary-layer energy: $\partial_t (\rho_{\mathrm{boundary}} V_{\mathrm{boundary}}) + \dots$
  \item Decay mechanisms: dissipation (viscosity, thermal conduction), dilution (expansion), leakage (transmission into condensate).
  \item Dimensional analysis and scaling relations: $\fmech \sim f(\Zpf/\Zcond, \uf/\cs, \ell_{\mathrm{boundary}}/H^{-1})$.
  \item Numerical integration of boundary-layer dynamics to produce $\umech(a)$ and verify consistency with Equation \eqref{eq:u_mech}.
\end{itemize}

\paragraph{Pass criterion.}
A physically motivated mapping from $(\Zpf, \Zcond, \uf, \ell_{\mathrm{boundary}})$ to $(\fmech, \ac, \sigmamech)$ that reproduces the best-fit phenomenology from Paper 1.

\MATHAGENT~\textbf{TODO:} Set up the boundary-layer dynamical equations (conservation laws, impedance matching, dissipation terms). Solve numerically for a range of initial conditions and verify the emergence of a Gaussian-like $\umech(a)$ profile.

\subsection{Paper 5: Structure Growth and Lensing}
\label{sec:paper5}

\paragraph{Objective.}
Verify that the condensate framework's preferred parameter region does not worsen the $S_8$ tension and remains consistent with large-scale structure observables: matter power spectrum $P(k)$, weak lensing shear $C_\ell^{\kappa\kappa}$, CMB lensing $C_\ell^{\phi\phi}$, BAO full-shape fits, redshift-space distortions (RSD), and cluster abundance $\sigma_8(z)$.

\paragraph{Key deliverables.}
\begin{itemize}[leftmargin=1.8em]
  \item Linear perturbation analysis: growth factor $D(a)$, growth rate $f(a) = \dx \ln D/\dx \ln a$, and their dependence on $\umech(a)$.
  \item Matter power spectrum $P(k)$ at $z = 0$ from CLASS/CAMB with modified background.
  \item Weak lensing convergence power spectrum $C_\ell^{\kappa\kappa}$ and comparison to DES, KiDS, HSC data.
  \item CMB lensing power spectrum $C_\ell^{\phi\phi}$ and comparison to Planck, ACT, SPT.
  \item Derived $\sigma_8$ and $S_8$ from best-fit cosmologies; comparison to Planck CMB-only and LSS-only values.
  \item Redshift-space distortions: $f\sigma_8(z)$ measurements from BOSS, eBOSS, DESI.
\end{itemize}

\paragraph{Pass criterion.}
$S_8$ from condensate framework lies within the current $\sim 2\sigma$ Planck--LSS discrepancy band; no new tensions introduced.

\MATHAGENT~\textbf{TODO:} Run CLASS/CAMB linear perturbation solver with modified background and compute all LSS observables. Interface with weak lensing likelihood codes (e.g., CosmoSIS, Cosmosis-Standard-Library).

\TODO{PHYS-AGENT: Update LSS references to DESI 2024 BAO+RSD results, DES Year 6, Euclid early release.}

\subsection{Paper 6: Planet Formation in a Condensate}
\label{sec:paper6}

\paragraph{Objective.}
Provide a micro-to-macro aggregation pathway for planet formation within the condensate framework. Address how the condensate's coherence length, stiffness, and quantum pressure affect the collapse and fragmentation of molecular clouds, the formation of protoplanetary disks, and the timescales for planetesimal accretion. Check consistency with observational benchmarks: Solar System chronology (CAI formation $\sim 4.567$ Gyr ago), exoplanet demographics, and protoplanetary disk lifetimes ($\sim 1$--$10$ Myr).

\paragraph{Key deliverables.}
\begin{itemize}[leftmargin=1.8em]
  \item Modified Jeans instability analysis in the condensate: critical mass $M_J \sim \cs^3 / \sqrt{G^3 \rho}$, with $\cs$ and $\rho$ determined by condensate properties.
  \item Fragmentation cascade from molecular cloud ($\sim 10^4 M_\odot$) to protostar ($\sim 1 M_\odot$) to planetesimals ($\sim 10^{-6} M_\odot$).
  \item Accretion timescales: runaway growth, oligarchic growth, giant-impact phase; comparison to N-body simulations.
  \item Role of condensate coherence length $\ell_{\mathrm{coh}}$: does it set a minimum fragment size or affect disk viscosity?
  \item Comparison to Solar System isotopic chronology (Pb-Pb, Hf-W, I-Xe dating) and meteorite evidence.
\end{itemize}

\paragraph{Pass criterion.}
The condensate framework reproduces Solar System formation timescales and exoplanet occurrence rates to within current observational uncertainties; no new fine-tuning required.

\MATHAGENT~\textbf{TODO:} Derive the modified Jeans mass and cooling timescales in the condensate. Compare to standard star-formation theory (e.g., Larson's laws, Kennicutt-Schmidt relation).

\TODO{PHYS-AGENT: Reference exoplanet demographics from Kepler, TESS, Gaia, and direct imaging surveys. Discuss whether condensate physics offers any new predictions for planet occurrence rates, metallicity dependence, or disk dispersal mechanisms.}

\subsection{Paper 7: Targeted Null Tests and Discriminants}
\label{sec:paper7}

\paragraph{Objective.}
Design observational measurements that must be null (or must deviate from $\Lambda$CDM) if the boundary-layer epoch occurred as predicted by the condensate framework. These ``smoking gun'' tests will either corroborate or falsify the framework.

\paragraph{Key deliverables.}
\begin{itemize}[leftmargin=1.8em]
  \item \textbf{Primordial power spectrum:} Does the boundary-layer bump imprint features in $P_{\mathrm{primordial}}(k)$ at specific scales corresponding to the horizon size at $a = \ac$? Forecast detectability with CMB-S4, LiteBIRD, 21-cm intensity mapping.
  \item \textbf{Gravitational-wave stochastic background:} Does the boundary epoch generate a relic stochastic GW background at frequencies $f \sim H(\ac) / (2\pi) \sim 10^{-17}$--$10^{-15}$ Hz? Forecast sensitivity of pulsar timing arrays (NANOGrav, SKA) and space-based detectors (LISA, BBO).
  \item \textbf{Non-Gaussianity:} Does the boundary-layer energy source non-Gaussian correlations in the CMB or LSS? Compute $f_{\mathrm{NL}}^{\mathrm{local}}$, $f_{\mathrm{NL}}^{\mathrm{equil}}$, $f_{\mathrm{NL}}^{\mathrm{ortho}}$ from boundary-layer perturbations.
  \item \textbf{Birefringence:} Could the condensate's coherence properties induce CMB polarization rotation (cosmic birefringence)? Recent Planck analyses report $\beta \sim 0.3^\circ$ rotation \cite{PlanckBirefringence2020}; is this explainable by boundary-layer physics?
  \item \textbf{Isocurvature modes:} Does the phase transition generate isocurvature (entropy) perturbations in addition to adiabatic curvature perturbations? Forecast constraints from CMB+BAO.
\end{itemize}

\paragraph{Pass criterion.}
At least one of these targeted tests yields a prediction that is distinguishable from $\Lambda$CDM and testable with near-future observatories (CMB-S4, Simons Observatory, LiteBIRD, LISA, SKA, Rubin LSST, Euclid, Roman).

\MATHAGENT~\textbf{TODO:} Calculate the primordial power spectrum modulation from the boundary-layer epoch. Compute the stochastic GW background spectrum. Evaluate bispectrum $\langle \zeta_{k_1} \zeta_{k_2} \zeta_{k_3} \rangle$ for primordial non-Gaussianity.

\TODO{PHYS-AGENT: Survey recent null-test literature (e.g., EDE smoking guns, modified gravity discriminants, early-Universe phase transitions and their GW signatures). Reference NANOGrav 15-year results, JWST high-$z$ galaxy anomalies, and any other emerging tensions that the condensate framework might address.}

\GRAPHAGENT~\textbf{Figure 10 placeholder:} Summary flowchart of the seven follow-on papers, showing dependencies and key deliverables for each.

% ============== SECTION 7: CONCLUSIONS ==============
\section{Conclusions}
\label{sec:conclusions}

We have presented a foundational framework in which the observable Universe is a coherent quantum condensate expanding into a universal protofluid. The Big Bang is reinterpreted as a phase-transition nucleation event, and cosmic expansion corresponds to the outward propagation of a conversion boundary. Gravitational response emerges from the condensate's effective stiffness---a mechanical-geometric linkage established in prior work \cite{SkavbergMechG2025,SkavbergMechGR2025,SkavbergPlanckField2025}---while boundary-layer impedance mismatch mediates transient energy reservoirs that can briefly modify the early expansion history.

This framework is \emph{falsifiable by design}. In equilibrium, it reduces exactly to Einstein--Hilbert dynamics, recovering all classical GR tests (Solar System, binary pulsars) and satisfying the GW170817 luminal bound on gravitational-wave propagation. Deviations from standard physics are permitted only in a narrow early-time window (post-BBN, pre-recombination), where a transient energy component $\umech(a)$ can briefly raise $H(a)$ in an EDE-like manner. This modification must pass stringent validation checkpoints: BBN light-element abundances, Planck CMB acoustic structure, BAO sound-horizon integrity, SNe distance moduli, distance-ladder $H_0$ measurements, and large-scale structure growth history.

We have formalized these checkpoints as quantitative pass/fail criteria (Section \ref{sec:validation}) and enumerated a sequenced research roadmap comprising seven follow-on papers (Section \ref{sec:roadmap}) to stress-test the framework empirically. The first paper will implement the phenomenological $\umech(a)$ bump in CMB Boltzmann codes and assess whether it can address the $H_0$ tension without degrading CMB+BAO+SNe fits. Subsequent papers will verify BBN compatibility, prove luminal GW propagation, derive the boundary-layer kinematics from first principles, analyze structure growth and weak lensing, explore planet formation in the condensate paradigm, and design targeted null tests to distinguish the framework from $\Lambda$CDM and competing alternatives.

If the condensate framework passes all these tests, it will constitute a major paradigm shift: spacetime geometry and gravitational dynamics are not fundamental, but emergent properties of an underlying condensed-matter-like substrate. Conversely, failure at any single checkpoint---deviation from PPN constraints, anomalous $c_T$, BBN inconsistency, CMB misfit, BAO/SNe degradation, or LSS tension---will falsify the framework. This clarity of prediction is the hallmark of a scientifically robust theory.

\subsection{Outlook}
\label{sec:outlook}

The condensate framework opens several exciting avenues for future research beyond the seven papers outlined in Section \ref{sec:roadmap}:

\paragraph{Quantum origins of the protofluid.}
What is the microphysical nature of the protofluid? Is it a vacuum-like state with quantum fluctuations (akin to the false vacuum in inflationary scenarios), or a more exotic quantum liquid (e.g., a spin network, loop quantum gravity foam, causal set)? Can we derive the protofluid's equation of state and impedance from first principles in quantum field theory or quantum gravity?

\paragraph{Connection to inflation.}
Does the condensate framework obviate the need for inflation, or is inflation still required to generate primordial density perturbations? Could the nucleation event itself be inflationary (e.g., via a first-order phase transition with supercooling), producing both the condensate and the initial power spectrum $P_{\mathrm{primordial}}(k)$?

\paragraph{Dark matter and dark energy.}
Could dark matter be excitations or defects in the condensate (e.g., solitons, vortices, phonons)? Could dark energy be the residual tension at the conversion boundary, or a relic of the protofluid's vacuum energy? Can the framework address the coincidence problem (why $\rho_{\Lambda} \sim \rho_m$ today)?

\paragraph{Black holes in a condensate.}
How does the condensate framework describe black hole formation, horizons, Hawking radiation, and information paradox? Is the singularity at $r = 0$ replaced by a condensate core, or does it remain a true singularity? Can the framework shed light on the firewall paradox or complementarity?

\paragraph{Experimental tests.}
Can analogue-gravity experiments (BEC in optical lattices, superfluid helium, water waves) simulate aspects of the condensate cosmology and test the boundary-layer kinematics in the lab? Can table-top experiments constrain the coherence length $\ell_{\mathrm{coh}}$ or the stiffness parameters of the cosmic condensate?

\TODO{PHYS-AGENT: Expand the outlook subsection with references to recent quantum-gravity phenomenology, black-hole thermodynamics in analogue systems, and experimental tests of emergent gravity.}

\paragraph{Acknowledgments.}
I thank colleagues working on analogue gravity, early-Universe cosmology, and cosmological inference for discussions that shaped this foundational framework. I am grateful to the Planck, BOSS, eBOSS, DES, SH0ES, and LIGO/Virgo/KAGRA collaborations for making their data publicly available, enabling rigorous validation of theoretical frameworks. This research made use of the CLASS and CAMB Boltzmann codes, the NASA Astrophysics Data System, and arXiv.

\TODO{PHYS-AGENT: Personalize acknowledgments once collaborators are identified. Acknowledge funding sources if applicable.}

% ============== REFERENCES ==============
\begin{thebibliography}{99}
\setlength{\itemsep}{2pt}

% --- Author's prior papers (foundational) ---
\bibitem{SkavbergMechG2025}
R.~Skavberg, ``A Mechanical Route to $G$,'' Zenodo (2025). doi:10.5281/zenodo.18496765.

\bibitem{SkavbergMechGR2025}
R.~Skavberg, ``A Mechanical Extension of General Relativity: The Singularity in Equilibrium,'' Zenodo (2025). doi:10.5281/zenodo.18393278.

\bibitem{SkavbergPlanckField2025}
R.~Skavberg, ``Planck Field --- Resolving Hulse--Taylor Binary,'' Zenodo (2025). doi:10.5281/zenodo.18462909.

% --- CMB and Planck ---
\bibitem{Planck2018}
Planck Collaboration (N.~Aghanim \emph{et al.}), ``Planck 2018 results. VI. Cosmological parameters,'' \emph{Astron.\ Astrophys.}\ \textbf{641}, A6 (2020). arXiv:1807.06209.

\bibitem{PlanckS8}
Planck Collaboration (N.~Aghanim \emph{et al.}), ``Planck 2018 results. VIII. Gravitational lensing,'' \emph{Astron.\ Astrophys.}\ \textbf{641}, A8 (2020). arXiv:1807.06210.

\bibitem{PlanckLensing}
Planck Collaboration (N.~Aghanim \emph{et al.}), ``Planck 2018 results. VIII. Gravitational lensing,'' \emph{Astron.\ Astrophys.}\ \textbf{641}, A8 (2020).

\bibitem{PlanckBirefringence2020}
Planck Collaboration (Y.~Minami \emph{et al.}), ``New Extraction of the Cosmic Birefringence from the Planck 2018 Polarization Data,'' \emph{Phys.\ Rev.\ Lett.}\ \textbf{125}, 221301 (2020). arXiv:2011.11254.

\bibitem{ACTLensing}
ACT Collaboration (M.~S.~Madhavacheril \emph{et al.}), ``The Atacama Cosmology Telescope: CMB lensing at $z\sim 0.5$,'' \emph{Astrophys.\ J.}\ \textbf{962}, 113 (2024). arXiv:2304.05203.

% --- Hubble tension and EDE ---
\bibitem{RiessReview2024}
A.~G.~Riess and L.~Breuval, ``The Local Value of $H_0$,'' arXiv:2308.10954 (2024).

\bibitem{Riess2022}
A.~G.~Riess \emph{et al.}, ``A Comprehensive Measurement of the Local Value of the Hubble Constant with 1 km s$^{-1}$ Mpc$^{-1}$ Uncertainty from the Hubble Space Telescope and the SH0ES Team,'' \emph{Astrophys.\ J.\ Lett.}\ \textbf{934}, L7 (2022). arXiv:2112.04510.

\bibitem{FreedmanTRGB2020}
W.~L.~Freedman \emph{et al.}, ``Calibration of the Tip of the Red Giant Branch (TRGB),'' \emph{Astrophys.\ J.}\ \textbf{891}, 57 (2020). arXiv:2002.01550.

\bibitem{JWST_Cepheids2023}
A.~G.~Riess \emph{et al.}, ``JWST Observations Reject Unrecognized Crowding of Cepheid Photometry as an Explanation for the Hubble Tension at $8\sigma$ Confidence,'' \emph{Astrophys.\ J.\ Lett.}\ \textbf{956}, L18 (2023). arXiv:2307.15806.

\bibitem{EDEReview2023}
M.~Kamionkowski and A.~G.~Riess, ``The Hubble Tension and Early Dark Energy,'' \emph{Annu.\ Rev.\ Nucl.\ Part.\ Sci.}\ \textbf{73}, 153 (2023). arXiv:2211.04492.

\bibitem{PoulinReview2023}
V.~Poulin, T.~L.~Smith, T.~Karwal, ``The Ups and Downs of Early Dark Energy solutions to the Hubble tension,'' \emph{Phys.\ Dark Univ.}\ \textbf{42}, 101348 (2023). arXiv:2302.09032.

\bibitem{HillEDE2020}
J.~C.~Hill, E.~McDonough, M.~W.~Toomey, S.~Alexander, ``Early Dark Energy Does Not Restore Cosmological Concordance,'' \emph{Phys.\ Rev.\ D}\ \textbf{102}, 043507 (2020). arXiv:2003.07355.

% --- BBN ---
\bibitem{BBN_Gbounds}
J.~Alvey, N.~Sabti, M.~Escudero, M.~Fairbairn, ``Improved BBN constraints on the variation of the gravitational constant,'' \emph{Eur.\ Phys.\ J.\ C}\ \textbf{80}, 148 (2020). arXiv:1910.10730.

\bibitem{PDG_BBN}
R.~L.~Workman \emph{et al.}\ (Particle Data Group), ``Review of Particle Physics,'' \emph{Prog.\ Theor.\ Exp.\ Phys.}\ \textbf{2022}, 083C01 (2022). [BBN section.]

\bibitem{Cooke_D2014}
R.~J.~Cooke, M.~Pettini, R.~A.~Jorgenson, M.~T.~Murphy, C.~C.~Steidel, ``Precision Measures of the Primordial Abundance of Deuterium,'' \emph{Astrophys.\ J.}\ \textbf{781}, 31 (2014). arXiv:1308.3240.

\bibitem{AdesOlivares_He2020}
E.~Aver, K.~A.~Olive, E.~D.~Skillman, ``The effects of He I $\lambda$10830 on helium abundance determinations,'' \emph{J.\ Cosmol.\ Astropart.\ Phys.}\ \textbf{07}, 003 (2015). arXiv:1503.08146.

% --- GW and GR tests ---
\bibitem{GW170817}
B.~P.~Abbott \emph{et al.}\ (LIGO Scientific \& Virgo Collaborations), ``Gravitational Waves and Gamma-Rays from a Binary Neutron Star Merger: GW170817 and GRB 170817A,'' \emph{Astrophys.\ J.\ Lett.}\ \textbf{848}, L13 (2017). arXiv:1710.05834.

\bibitem{PPN_Cassini}
B.~Bertotti, L.~Iess, P.~Tortora, ``A test of general relativity using radio links with the Cassini spacecraft,'' \emph{Nature}\ \textbf{425}, 374 (2003).

\bibitem{PPN_Perihelion}
E.~V.~Pitjeva, N.~P.~Pitjev, ``Relativistic effects and dark matter in the Solar system from observations of planets and spacecraft,'' \emph{Mon.\ Not.\ R.\ Astron.\ Soc.}\ \textbf{432}, 3431 (2013). arXiv:1306.3043.

\bibitem{PPN_Shapiro}
B.~Bertotti, L.~Iess, P.~Tortora, ``A test of general relativity using radio links with the Cassini spacecraft,'' \emph{Nature}\ \textbf{425}, 374 (2003).

\bibitem{HulseTaylor}
J.~M.~Weisberg, D.~J.~Nice, J.~H.~Taylor, ``Timing Measurements of the Relativistic Binary Pulsar PSR B1913+16,'' \emph{Astrophys.\ J.}\ \textbf{722}, 1030 (2010). arXiv:1011.0718.

\bibitem{DoublePulsar}
M.~Kramer \emph{et al.}, ``Tests of General Relativity from Timing the Double Pulsar,'' \emph{Science}\ \textbf{314}, 97 (2006). arXiv:astro-ph/0609417.

% --- BAO and SNe ---
\bibitem{BOSS_BAO}
S.~Alam \emph{et al.}\ (BOSS Collaboration), ``The clustering of galaxies in the completed SDSS-III Baryon Oscillation Spectroscopic Survey: cosmological analysis of the DR12 galaxy sample,'' \emph{Mon.\ Not.\ R.\ Astron.\ Soc.}\ \textbf{470}, 2617 (2017). arXiv:1607.03155.

\bibitem{eBOSS_BAO}
S.~Alam \emph{et al.}\ (eBOSS Collaboration), ``Completed SDSS-IV extended Baryon Oscillation Spectroscopic Survey: Cosmological implications from two decades of spectroscopic surveys at the Apache Point Observatory,'' \emph{Phys.\ Rev.\ D}\ \textbf{103}, 083533 (2021). arXiv:2007.08991.

\bibitem{Pantheon}
D.~M.~Scolnic \emph{et al.}, ``The Complete Light-curve Sample of Spectroscopically Confirmed SNe Ia from Pan-STARRS1 and Cosmological Constraints from the Combined Pantheon Sample,'' \emph{Astrophys.\ J.}\ \textbf{859}, 101 (2018). arXiv:1710.00845.

\bibitem{PantheonPlus}
D.~Brout \emph{et al.}\ (Pantheon+ Collaboration), ``The Pantheon+ Analysis: Cosmological Constraints,'' \emph{Astrophys.\ J.}\ \textbf{938}, 110 (2022). arXiv:2202.04077.

% --- LSS and weak lensing ---
\bibitem{DES2021}
DES Collaboration (T.~M.~C.~Abbott \emph{et al.}), ``Dark Energy Survey Year 3 results: Cosmological constraints from galaxy clustering and weak lensing,'' \emph{Phys.\ Rev.\ D}\ \textbf{105}, 023520 (2022). arXiv:2105.13549.

\bibitem{KiDS2020}
KiDS Collaboration (M.~Asgari \emph{et al.}), ``KiDS-1000 Cosmology: Cosmic shear constraints and comparison between two point statistics,'' \emph{Astron.\ Astrophys.}\ \textbf{645}, A104 (2021). arXiv:2007.15633.

\bibitem{LymanAlpha}
J.~E.~Bautista \emph{et al.}, ``The completed SDSS-IV extended Baryon Oscillation Spectroscopic Survey: measurement of the BAO and growth rate of structure of the luminous red galaxy sample from the anisotropic correlation function between redshifts 0.6 and 1.0,'' \emph{Mon.\ Not.\ R.\ Astron.\ Soc.}\ \textbf{500}, 736 (2021). arXiv:2007.08993.

% --- Induced gravity and analogue gravity ---
\bibitem{Sakharov}
A.~D.~Sakharov, ``Vacuum quantum fluctuations in curved space and the theory of gravitation,'' \emph{Sov.\ Phys.\ Dokl.}\ \textbf{12}, 1040 (1968). [Reprinted in \emph{Gen.\ Rel.\ Grav.}\ \textbf{32}, 365 (2000).]

\bibitem{Vassilevich}
D.~V.~Vassilevich, ``Heat kernel expansion: User's manual,'' \emph{Phys.\ Rept.}\ \textbf{388}, 279 (2003). arXiv:hep-th/0306138.

\bibitem{AnalogueGravityReview}
C.~Barceló, S.~Liberati, M.~Visser, ``Analogue Gravity,'' \emph{Living Rev.\ Relativity}\ \textbf{14}, 3 (2011). arXiv:gr-qc/0505065.

\end{thebibliography}

% ============== APPENDICES (placeholders) ==============
\newpage
\appendix

\section{Derivation of Rankine--Hugoniot Relations with Surface Stress-Energy}
\label{app:RH_derivation}

\MATHAGENT~\textbf{TODO:} Provide detailed step-by-step derivation of Equation \eqref{eq:RH_general} from stress-energy conservation, Gauss--Stokes theorem, and matching conditions at $\Sigma$.

\section{Impedance Matching and Reflection Coefficients}
\label{app:impedance}

\MATHAGENT~\textbf{TODO:} Derive Equation \eqref{eq:reflection_transmission} from linearized wave equations at the boundary, matching conditions for pressure and velocity, and energy flux conservation.

\section{Calculation of Sound Horizon with Transient Energy Component}
\label{app:sound_horizon}

\MATHAGENT~\textbf{TODO:} Integrate $\rs = \int_0^{a_{\mathrm{dec}}} \cs(a)/[a^2 H(a)] \, \dx a$ with the modified Friedmann equation \eqref{eq:Friedmann_general}. Show how $\umech(a)$ reduces $\rs$ by a few percent for representative $(\fmech, \ac, \sigmamech)$.

\section{Tensor-Mode Propagation in Modified FLRW Background}
\label{app:tensor_modes}

\MATHAGENT~\textbf{TODO:} Linearize the Einstein equations for tensor perturbations $h_{ij}$ on the FLRW background with $\umech(a)$. Show that $c_T = c$ at late times and compute any transient friction or damping during the boundary epoch.

\section{Glossary of Notation}
\label{app:notation}

\TODO{PHYS-AGENT: Compile a comprehensive glossary of all symbols, subscripts, and abbreviations used in the paper. Include definitions for $\Zpf$, $\Zcond$, $\umech$, $\fmech$, $\ac$, $\sigmamech$, $\rs$, $\cs$, $\mpl$, etc.}

\end{document}
